\documentclass[11pt,xcolor=table,svgnames,aspectratio=169]{beamer}
\usepackage{graphicx}
\usepackage{fontspec}
\usepackage{tikz}
\usepackage{doi}
\usepackage{fontawesome}
\usepackage{academicons}
\usepackage{xcolor}
\usepackage{tabularx}
\usepackage{tcolorbox}
\usepackage{booktabs}
\usepackage{colortbl}
\usepackage{stackengine}
\usepackage{bbding}


%--------------------------------
\hypersetup{bookmarksopen=true,
bookmarksnumbered=true,  
pdffitwindow=true, 
pdfstartview=Fit,
pdffitwindow=true,
pdftoolbar=true,
pdfmenubar=true,
pdfwindowui=true,
pdfsubject={EMODnet Biology, Data analysis, Interpolation, DIVAnd},
pdfauthor={C. Troupin},
bookmarksopenlevel=0,
colorlinks=true,
linkcolor=frametitle,anchorcolor=black,%
citecolor=frametitle,filecolor=black,%
menucolor=black,urlcolor=frametitle,%
pdfpageduration=1,%
}


\newcolumntype{Q}{>{\raggedleft\arraybackslash}X} 
\newcolumntype{P}{>{\raggedright\arraybackslash}X}

\usepackage{array}
\newcolumntype{L}[1]{>{\raggedright\let\newline\\\arraybackslash\hspace{0pt}}m{#1}}
\newcolumntype{C}[1]{>{\centering\let\newline\\\arraybackslash\hspace{0pt}}m{#1}}
\newcolumntype{R}[1]{>{\raggedleft\let\newline\\\arraybackslash\hspace{0pt}}m{#1}}


%\usefonttheme{serif}
\usefonttheme{professionalfonts}
\usetikzlibrary{positioning}

\DeclareGraphicsExtensions{.pdf,.png,.PNG,.JPG,.jpg,.jpeg,.gif}
\graphicspath{
{./figures/},
{/home/ctroupin/Presentations/figures4presentations/},
{/home/ctroupin/Presentations/figures4presentations/logo/},
}

\usetikzlibrary{arrows,shapes,backgrounds,spy,calc}
\tikzstyle{every picture}+=[remember picture]
\tikzstyle{na} = [baseline=-.5ex]

\newcommand*\circled[4]{\tikz[baseline=(char.base)]{
    \node[shape=circle, fill=#2, draw=#3, text=#4, inner sep=2pt, opacity=.75] (char) {#1};}}

\usepackage{dashbox}
\newcommand{\mybox}[1]{\color{blueemodnet}{\dbox{#1}}}



\urlstyle{sf}

\newcommand{\logoheight}{.75cm}

\setmainfont{Carlito}

\definecolor{author}{HTML}{7F7F7F}
\definecolor{frametitle}{HTML}{333333}
\definecolor{email}{HTML}{FFC000}
\definecolor{blueemodnet}{HTML}{0a71b4}


\setbeamercolor{title}{fg=white}
\setbeamercolor{author}{fg=white}
\setbeamercolor{institute}{fg=white}
\setbeamercolor{frametitle}{fg=frametitle}
\setbeamercolor{date}{fg=white}
\setbeamercolor{enumerate item}{fg=blueemodnet}
\setbeamercolor{description item}{fg=blueemodnet}
\setbeamercolor{itemize item}{fg=blueemodnet}


\setbeamerfont{title}{size=\fontsize{14pt}{10}\selectfont}
\setbeamerfont{author}{size=\fontsize{9.5pt}{5}\selectfont}
\setbeamerfont{subtitle}{size=\fontsize{8pt}{10}\selectfont}
\setbeamerfont{frametitle}{size=\fontsize{12pt}{8}\selectfont,series=\bfseries}
\setbeamerfont{date}{size=\fontsize{8pt}{10}\selectfont}
\setbeamerfont{page number in head/foot}{size=\tiny}

\setbeamertemplate{frametitle}{\vspace*{.3cm}\hspace*{-0.5cm}\insertframetitle}
%\defbeamertemplate{itemize item}{image}{\small\includegraphics[height=1.6ex]{itememodnet.png}}
\setbeamertemplate{itemize item}[image]
\setbeamertemplate{navigation symbols}{}

%\setbeamertemplate{frametitle}{\hspace*{-0.5cm}\insertframetitle}

\setbeamertemplate{footline}
{
  \hbox{%
  \begin{beamercolorbox}[wd=.97\paperwidth,ht=2.25ex,dp=1ex,right]{date in head/foot}%
    \usebeamerfont{footline} \textcolor{frametitle}{\insertframenumber{}}
  \end{beamercolorbox}}%
  \vskip0.35cm%
}

\definecolor{commentcolor}{rgb}{0.75,.75,.75}
\renewcommand{\comment}[1]{\hfill\textcolor{commentcolor}{\footnotesize{#1}}}


\usepackage[export]{adjustbox}
\newcommand{\listcheck}{\textcolor{DarkGreen}{\checkmark}}
\newcommand{\listcross}{\textcolor{red}{\XSolidBrush}}

\title{EMODnet Biology Phase VI}
\subtitle{Progress meeting}
\author{C.~Troupin}
\institute{University of Liège}
\date{12-13 May 2025}
\titlegraphic{\includegraphics[height=\logoheight]{logo_uliege}~\includegraphics[height=\logoheight]{logo_gher}
}

\definecolor{highlightext}{HTML}{8E8E8E}

\AtBeginDocument{\usebeamerfont{normal text}}

\begin{document}


% Title frame, no content needed here
{
\thispagestyle{empty}
\usebackgroundtemplate{\includegraphics[width=\paperwidth]{pg_0001.pdf}}
\begin{frame}[noframenumbering]
\end{frame}
}
% -----------------------------------
\usebackgroundtemplate{\includegraphics[width=\paperwidth]{pg_0002.pdf}}

\begin{frame}
\frametitle{Graphical abstract}

How to we go from this \ldots \hspace*{4cm} \onslide<2>{to this?}\\
(sparse osbervations) \hspace*{4cm} \onslide<2>{(monthly gridded fields)}\\
\begin{figure}
\flushleft
\includegraphics<1->[width=.47\textwidth]{Small_copepods_raw_201604.png}\includegraphics<2>[width=.47\textwidth]{Small_copepods_best_201604.png}
\end{figure}

\end{frame}

% -------------------------------------------------------------------------------------------------------

\begin{frame}
\frametitle{The results: 2 netCDF files storing monthly gridded fields}

\begin{columns}

\begin{column}{0.45\textwidth}

\begin{figure}
\includegraphics[width=\textwidth]{Small_copepods_best_196309.png}
\end{figure}

\end{column}

\begin{column}{0.55\textwidth}

\begin{description}
\footnotesize
\item[Variables:] abundances of specific taxonomic groups:\\
1.~small copepods \&  2.~large copepods
\item[Spatial resolution:] 1° $\times$ 1°
\item[Temporal resolution:] monthly fields
\item[Temporal coverage:] 1958 -- 2022
\item[Download:] ~ \\ \url{https://dox.uliege.be/index.php/s/DkhXlKVQpDXVWaE/download}
\end{description}
  
\end{column}


\end{columns}


\end{frame}

% -------------------------------------------------------------------------------------------------------

\begin{frame}
\frametitle{The data: Continuous Plankton Recorder (CPR) survey}

\begin{columns}

\begin{column}{0.45\textwidth}

\begin{figure}
\includegraphics<1>[width=.9\textwidth]{data_location_NortheastAtlantic.png}
\includegraphics<2>[width=\textwidth]{numobs_grid_Small_copepods.png}
\includegraphics<3>[width=\textwidth]{Small_copepods_raw_198106.png}

\end{figure}

\end{column}

\begin{column}{0.55\textwidth}

\begin{itemize}
\item Long time series with consistent measurement technique (since 1931)
\item Extensive spatial coverage
\item[]
\item[\aiDoi] \href{https://doi.org/10.17031/68da4a97650f1}{10.17031/68da4a97650f1}
\item[\faLink] \url{https://doi.mba.ac.uk/data/3548/}
  
\end{itemize}

\end{column}
\end{columns}
\end{frame}

% -------------------------------------------------------------------------------------------------------


\begin{frame}
\frametitle{The gridding tool: DINCAE.jl}

\begin{figure}
\centering
\includegraphics[height=1cm]{logo_julia}~~\includegraphics[height=1cm]{logo_gpu}~~\includegraphics[height=1cm]{logo_jupyter.png}
\end{figure}

Data INterpolating Convolutional Auto-Encoder (DINCAE)
\begin{description}
\item[\faGithub] \url{https://github.com/gher-uliege/DINCAE.jl}
\item[\aiDoi] \href{https://doi.org/10.5281/zenodo.5575066}{10.5281/zenodo.5575066}
\item[\faFilePdfO] \url{https://gmd.copernicus.org/articles/15/2183/2022/}
\end{description}


\end{frame}

% -------------------------------------------------------------------------------------------------------


\begin{frame}
\frametitle{Specific developments of DINCAE}

\begin{columns}
\begin{column}{0.5\textwidth}

\begin{enumerate}
\item Use \textit{environmental variables} to help \textit{guide} the neural network
\begin{itemize}
\item Sea surface temperature
\item Bathymetry
\item Distance to coast
\item \ldots
\end{itemize}
\item Add \textit{smoothing constrain} to the cost function 
\end{enumerate}



\end{column}

\begin{column}{0.5\textwidth}
\begin{figure}
\centering
\includegraphics[width=.95\textwidth]{hadley_sst.png}
\end{figure}

\end{column}
\end{columns}




\end{frame}


% -------------------------------------------------------------------------------------------------------

\begin{frame}
\frametitle{Sequence of operations}


\begin{enumerate}
\item<1-> Create training and validation dataset \comment{10\% of observations}
\item<2-> Generate random set of parameters \comment{number of epochs, learning rate, smoothing, \ldots}
\item<3-> Perform analysis \comment{with $\ne$ environment variables}
\item<4-> Compute Root Mean Square Deviation \comment{using the validation dataset}
\item<5-> Repeat steps 2 $\rightarrow$ 4 many times \comment{(100's)}
\item<6-> Find best combination of parameters \comment{i.e. with lowest RMSD}
\item<7-> Perform gridding with best parameters on full datasets 
\end{enumerate}


\end{frame}

% -------------------------------------------------------------------------------------------------------


\begin{frame}
\frametitle{Results validation}


\begin{columns}
\begin{column}{0.5\textwidth}
   \begin{enumerate}
\item<1-> Time series of spatially averaged values
\item<3-> Time-averaged monthly fields (June)
\item<4-> Time-averaged monthly fields (December)
\end{enumerate}
\end{column}
\begin{column}{0.5\textwidth}
    \begin{figure}
    \centering
     \includegraphics<1>[width=\textwidth]{spatial_avareged_Small_copepods_best.png} 
     \includegraphics<2>[width=\textwidth]{spatial_avareged_Small_copepods_best_2010_2022.png}   
     \includegraphics<3>[width=0.5\textwidth]{time_averaged_Small_copepods_best_06.png}~\includegraphics<3>[width=0.5\textwidth]{time_averaged_Small_copepods_rawdata_06.png}
     \includegraphics<4>[width=0.5\textwidth]{time_averaged_Small_copepods_best_12.png}~\includegraphics<4>[width=0.5\textwidth]{time_averaged_Small_copepods_rawdata_12.png}
     \end{figure}
\end{column}
\end{columns}

\end{frame}


% -------------------------------------------------------------------------------------------------------

\begin{frame}[t]
\frametitle{Conclusions \& lessons learned}


\begin{columns}[t]
\begin{column}{0.33\textwidth}
{\centering \huge \circled{1}{blueemodnet}{black}{white} }

\vspace*{1cm}

Neural network algorithm capable to generate smooth, consistent gridded fields from very sparse observations
\end{column}

\begin{column}{.33\textwidth}
{\centering \huge \circled{2}{blueemodnet}{black}{white} }

\vspace*{1cm}

Essential to optimise the hyperparameter values 

\end{column}
\begin{column}{0.33\textwidth}
{\centering \huge \circled{3}{blueemodnet}{black}{white} }

\vspace*{1cm}

Requirement to run a large number of experiments before the final results \faPlus~GPU computing resources

\end{column}
\end{columns}

\vfill

\faGithub~\url{https://github.com/gher-uliege/EMODnet-Biology-Interpolation-NN-Copepods}


\end{frame}

% -------------------------------------------------------------------------------------------------------

% Final slide, nothing to touch
{
\thispagestyle{empty}
\usebackgroundtemplate{\includegraphics[width=\paperwidth]{pg_0005.pdf}}
\begin{frame}

\end{frame}
}




\end{document}
